% Q4 integrated from zhh core v3 at ee23b200. Numbers: outputs/Q4/results.json.
\section{问题四：能力演进、技术贡献与算力放缓}
\label{sec:q4}

\subsection{能力口径与数据筛选}
榜单与资源元数据的公开出处分别见 Open LLM Leaderboard 2\cite{openllmv2} 和 Epoch AI Models Dataset\cite{epochmodels}；本文数值取题面附件截面。
令模型的能力 $S_i$ 为 C1/C2 六项 Benchmark 百分制得分的等权均值，仅对六项均有数值的记录定义。时间采用榜单提交日，参数量 $N$ 单位为十亿参数。开放性由 Epoch 权重标志及模型仓库许可证联合筛选；明确的 Epoch 不开放记录不因许可证代理重新纳入。主分析按基础预训练与后训练（chat、领域微调）分组；持续预训练、合并、多模态仅用于口径敏感性。许可证只能代理开放可用性，不能证明权重当前可下载。C2 去重后的扩展开放集合有 2616 个模型，主基础组 204 个，主后训练组 1620 个。

C8 对 1860 个可解析模型目录逐任务聚合 24 项 BBH，并登记 4 个损坏 JSON；均分会掩盖任务间的低分尾部。C7 的\texttt{max\_position\_embeddings}提供上下文长度的观测情景，供第三问的注意力成本计算。C3 包含 4599 行混合来源历史表，其中 79 个 Model/Year/Source 重复键；按来源与年份逐任务分析，不把历史报告中的\texttt{Average}直接拼接为同口径能力曲线。26 行历史报告中\texttt{Average}与六任务均值的平均绝对差为 9.101 分。

C4 对 3523 行资源元数据逐行审计。正有限 $N,C,D$ 的完整记录 850 行；主开放语言模型口径 81 行、放宽 $C/(6ND)$ 门禁后 93 行。这里 $C$ 为 FLOPs，$D$ 为十亿 token；$6ND$ 仅用于单位与来源一致性筛查，部分 Epoch 算力可能本身由 $N,D$ 推算，故比值接近 1 不构成独立机制验证。2021--2024 完整年度的主资源集合，其年度算力对数增长率为 $g_C=0.690315$，放宽集合为 $0.600552$；2025 年尚不完整，不纳入增长率拟合。

\subsection{同支持贡献的条件分解}
分别在同类型、共同参数格上比较早末两个自然月的六任务均分。设格 $j$ 在时点 $t$ 的权重、均值为 $w_{jt},m_{jt}$，先将两期权重对称标准化，定义
\[
\Delta_{\mathrm{scale}}=\frac12\sum_j (w_{j1}-w_{j0})(m_{j1}+m_{j0}),\qquad
\Delta_{\mathrm{within}}=\frac12\sum_j (w_{j1}+w_{j0})(m_{j1}-m_{j0}).
\]
两项严格相加为共同支持的标准化均值变化。全样本均值变化与标准化变化之差另列为组成及支持差，不能归入技术。
\begin{table}[htbp]\centering\small
\caption{主两月窗口、0.5 decade 参数格的条件分解（分）}
\begin{tabular}{lrrrr}\toprule
类型 & 全样本变化 & 参数分布项 & 格内时间项 & 组成/支持差\\\midrule
后训练 & 4.222 & -0.758 & 6.521 & -1.542\\
基础预训练 & -1.192 & -3.756 & 4.636 & -2.072\\
\bottomrule\end{tabular}
\end{table}
后训练标准化变化为 $5.763=-0.758+6.521$ 分，相应带符号份额为 $-13.15\%$ 与 $113.15\%$，分母是标准化变化而非全样本的 4.222 分。为回答题目的“规模扩张/非规模技术进步”两类贡献，本文采用\emph{条件工程约定}：同类型、同一参数格内的 $D$ 分布与训练资源结构在两月间近似稳定，跨格迁移的 $N$ 同时代理与其伴随的 $D$ 和算力扩张；于是参数分布项称为\emph{可见规模代理贡献}，格内时间项称为\emph{非规模技术代理贡献}。这个约定把未逐模型观测的 $D$ 归入随 $N$ 扩张的规模代理，而不是假装已单独估计 $D$ 的份额。若格内 $D$、后训练投入、模型选择或评测口径发生变化，6.521 分还包含这些因素，不能解释为纯技术因果效应。

同为后训练、两月窗口时，参数格宽为 $0.25/0.50/0.75$ decade 的可见规模代理份额分别为 $-16.47\%/-13.15\%/-50.08\%$，非规模剩余代理为 $116.47\%/113.15\%/150.08\%$；每个比例都以其自身共同支持标准化变化为分母。方向一致而幅度敏感，故主口径的结论是“在上述稳定性假设下，近期增益主要由非规模项解释，规模分布迁移为负向抵消”，并非普适技术占比。基础组的标准化变化仅 0.880 分，比例对分母过敏，不与后训练合并为总体份额。增加开发者控制后共同支持仅剩早期 6、末期 8 条，约占各期 $1.8\%$，不能据此推广总体。格宽、窗口、开放口径与开发者块抽样的其余敏感性见贡献表。图\ref{fig:q4-contribution}展示精确账面分解。
\begin{figure}[htbp]\centering
\includegraphics[width=.82\textwidth]{Q4/01_contributions.png}
\caption{能力变化的参数分布、格内时间和组成差；格内时间不是已识别的纯技术效应。}\label{fig:q4-contribution}
\end{figure}

\subsection{两月稳健前沿与放缓情景}
本节把同类型最近两个月提交模型的六项均分 $q_{0.9}$ 定义为\emph{稳健前沿代理}：它描述高分模型群的上沿典型水平，抑制单个极端榜单记录，并不等于未来最高模型的绝对上界。准备表中，基础组最近两月 $q_{0.9}=21.032$、同期最高 32.436、全期最高 38.441；后训练组对应为 40.445、47.217、51.231 分。因此基础组的 21.032 常数预测绝不能读成“未来最强基础模型只有 21.032 分”。预测对象保持 $q_{0.9}$，最高值只作边界口径对照；最高值自身的未来轨迹未由四窗口回测校准。
令 $z=\operatorname{logit}(S/100)$。比较常数基线、直接两月前沿趋势、均值资源动力学及直接 $q_{0.9}$ 分位资源动力学。后者在同类型训练窗口最小化分位 check loss，分位系数由线性规划求得；滚动训练只见截点前提交的模型版本，测试为之后两个完整自然月，C4 资源也按当时日期截断。四个测试窗口彼此重叠且已用于模型选择，是诊断回测而非盲测。后训练组四窗 RMSE 依次为 3.737、3.053、2.237、2.199 分，分位资源候选仅略优于均值候选；基础组常数 RMSE 7.633，优于分位资源候选 10.561。各模型 $R^2$ 仍为负，不能声称长期预测已经校准。

后训练组观测截点为 2025-03-13，基础组为 2025-03-09；以下 12/24 月均从\emph{各自附件截点}起算。后训练组的目标日为 2026-03-13 / 2027-03-13，基础组为 2026-03-09 / 2027-03-09，并非以本次竞赛日期为预测起点。设未来 $H$ 月算力对数增长为 $s g_C H/12$，其中 $s\in\{1,1/2,1/4,0\}$ 表示历史、减半、四分之一和停止。主情景假设 $\eta=1/2$ 的增量流向 $\log_{10}N$，其余流向 $\log_{10}D$；时间趋势保留系数为 $\lambda=1/2$。参数与 token 配置仅是资源情景，C2 没有逐模型 $D$，不能据此估计独立的 $D$ 能力效应。条件稳健前沿由近期两月 $q_{0.9}$ 锚点加上资源及时间项得到：
\[
\widehat S_k(H)=100\,\sigma\!\left[\operatorname{logit}(S_{k,0}/100)
 +\beta_{N,k}\eta s g_C H/12+\lambda\beta_{t,k}H\right].
\]
将类型内局部分位回归锚到近期总体 $q_{0.9}$ 是明示平移假设。

后训练组算力增速减半时，2026-03-13 的 12 月主候选为 53.828 分，2027-03-13 的 24 月为 66.681 分；固定主假设下 80 次块抽样的 12 月 5--95\% 诊断范围为 50.302--56.007。将模型、资源速度、时间保留、窗口和分位假设合并后，12 月条件范围为 37.755--72.236；再额外假设未来固定偏差不超过历史最大候选误差，则压力范围为 31.590--78.400。两者都不是具有已验证覆盖率的未来预测区间。基础组在该回测中以常数更优，2026-03-09 的 12 月稳健前沿参考值为 21.032；其资源响应曲线仅保留为压力模型。图\ref{fig:q4-frontier}给出回测与情景。
\begin{figure}[htbp]\centering
\includegraphics[width=.84\textwidth]{Q4/02_frontier_and_validation.png}
\caption{两类模型的局部前沿、四窗口回测及条件情景；长期外推没有覆盖率保证。}\label{fig:q4-frontier}
\end{figure}

\subsection{Loss--能力映射与上游误差}
C6 的 Loss 来自不同语料、分词器和训练阶段，不能混成单一横轴。按\texttt{Loss\_Source}及 validation/training/unspecified 标签分层，仅在来源内拟合单调候选并与常数基线留一比较。标为 Qwen2 validation 的 4 条记录 RMSE 为 7.048 对 10.708，Qwen2.5 的 6 条为 2.655 对 15.074；Pythia 高可比组却未优于常数。附件标签尚不能证明 Loss 与第三问预测值处于同一数值坐标。正式跨坐标接口因此保持\texttt{ready\_for\_empirical\_Q3\_conversion=false}。

为说明误差影响，在第三问冻结的 36 条固定配方政策格上，对 5 个来源分别施加 $L_{\rm source}=aL_{\rm Q3}+b$ 的明示情景，并对上游 Loss 加 $\pm0.05$ 压力；总计 4860 行。$a\in\{0.9,1,1.1\}$、$b\in\{-0.1,0,0.1\}$ 是假设网格，非数据估计的置信区间。每行单独保留上游不可行、参数超域或 Loss 超域；仅在来源支持域内比较同 $N$ 条件参考 Loss 与政策 Loss 的能力差。表\ref{tab:q4-q3-conversion} 选取第三问在中、高预算实际推荐的配方 172、8192 Token、幂成本，展示可直接阅读的条件换算。

\begin{table}[htbp]\centering\small
\caption{Q3 配方 172 的条件能力换算；分数范围是 $a,b,\Delta L$ 假设网格中可换算行的极值，不是置信区间}
\label{tab:q4-q3-conversion}
\begin{tabular}{llrrrr}\toprule
预算/FLOPs & C6 来源 & Q3 Loss & 中心能力分 & 可换算分数范围 & 相对同 $N$ 骨架增益范围\\\midrule
$10^{22}$ & Qwen2 validation & 2.0944 & 22.371 & 13.601--27.462 & 0.322--3.596\\
$10^{22}$ & Qwen2.5 validation & 2.0944 & 28.695 & 15.057--39.197 & 0.473--5.734\\
$10^{24}$ & Qwen2 validation & 2.0195 & 24.602 & 15.288--27.873 & 0.555--3.780\\
$10^{24}$ & Qwen2.5 validation & 2.0195 & 32.262 & 17.583--40.404 & 0.839--6.467\\
\bottomrule\end{tabular}
\end{table}
这里的中心是 $a=1,b=0,\Delta L=0$，同 $N$ 骨架仅为模型内参考，绝非观测 Benchmark 基线。上述两来源、两预算的增益在各自\emph{可换算}情景中均保持正号；其他来源或格可能变号，超域情景没有得分结论。对于同一来源、共同支持且 $a>0$ 的候选，来源内单调桥接保持 Q3 Loss 的次序，因此共同坐标误差改变得分幅度而不反转该组候选的次序；一旦候选超域，不能继续据此比较。此表消费中高预算的主推荐，低预算联立配方 477 与独立 $Q_B$ 模式尚不在这 36 条桥接格中，不能声称其能力排名已经换算。直接用 C2/C4 预测的稳健前沿不读取第三问 Loss，故桥接误差影响跨题政策解释，不改变独立前沿；两条路径共同支持的是\emph{支持域内的条件配置和情景能力判断}，而非同一个已校准未来分数。图\ref{fig:q4-bridge}展示来源内诊断。
\begin{figure}[htbp]\centering
\includegraphics[width=.84\textwidth]{Q4/03_loss_bridge.png}
\caption{来源内 Loss--能力候选及留一诊断；并非已标定的第三问跨坐标映射。}\label{fig:q4-bridge}
\end{figure}

\subsection{第四问结论及适用范围}
在同类型共同支持与“格内 $D$ 稳定、跨格 $N$ 代理规模投入”的工程约定下，近期后训练能力变化的可见规模代理贡献为 $-13.15\%$，非规模剩余代理为 $113.15\%$，格宽敏感性分别为 $-16.47\%$ 至 $-50.08\%$ 和 $116.47\%$ 至 $150.08\%$。算力减半情景给出从 2025-03-13 起算的后训练 12/24 月稳健前沿 53.828/66.681 分；它不表示未来最高模型上界。中高预算配方 172 的 Qwen 来源条件换算给出正增益，但分数范围较宽。数据没有识别纯技术因果贡献、第三问到 C6 的实证 Loss 坐标或长期预测覆盖率；这些限制在接口和结果表中保持原状态。
